
\subsection{Historical context and inherited methods}

Karatsuba's recursive construction, published jointly
with Ofman in 1962, gave a binary multiplication circuit of size
$O(n^{\log_2 3})$~\cite{KaratsubaOfman1962}. Toom's subsequent use of polynomial
evaluation and interpolation gave circuit size
$n\exp(O(\sqrt{\log n}))$, and hence $O(n^{1+\eta})$ for every fixed
$\eta>0$~\cite{Toom1963}. These early results were formulated in terms of
automata and logical networks. Cook adapted Toom's method to the
Turing-machine setting; the contemporary account of Sch\"onhage and
Strassen credits this adaptation explicitly
\cite[Section~1, p.~282]{SchonhageStrassen1971}.

In the multitape Turing model,
Sch\"onhage and Strassen proved the bound
$O(n\log n\log\log n)$ in 1971~\cite{SchonhageStrassen1971}.
F\"urer improved it to $n\log n\,2^{O(\log^* n)}$ in 2007, with a full
journal account in 2009~\cite{Furer2007,Furer2009}; here $\log^* n$ counts
iterated logarithms until a constant threshold is reached.
De, Kurur, Saha and Saptharishi developed a modular-arithmetic route to
the same scale~\cite{DeKururSahaSaptharishi2013}.
Later improvements made the exponential factor $8^{\log^* n}$ and then
$4^{\log^* n}$~\cite{HarveyHoevenLecerf2016,HarveyHoeven2019Lattice}.
Harvey and van der Hoeven proved the unconditional
$O(n\log n)$ bound~\cite{HarveyHoeven2021}.

The multitape bounds just cited concern that fixed model. Sch\"onhage obtained
linear-time multiplication on storage-modification machines~\cite{Schonhage1980}.
Groff later gave a smaller-than-$\log n$ factor in a unit-cost RAM model
with constant-time arbitrary memory access, after polynomial preprocessing
whose cost is excluded from that bound~\cite[Sections~2 and~4]{Groff2019}.
Neither result establishes the bound in the fixed finite-tape model of
Theorem~\ref{thm:main}.

The transform methods behind these bounds also underlie the present
construction. Sch\"onhage--Strassen multiplication uses roots of unity in
Fermat-type residue rings, where certain multiplications become shifts.
F\"urer combines cheap polynomial roots of unity with more general roots,
reducing the frequency of expensive multiplications. Harvey and van der
Hoeven use Gaussian resampling to pass to multidimensional transforms with
power-of-two dimensions, then evaluate them using polynomial transforms
developed by Nussbaumer and Quandalle and subsequently by
Nussbaumer~\cite{NussbaumerQuandalle1978,Nussbaumer1980,HarveyHoeven2021}.
The Gaussian-gridding work of Dutt and Rokhlin is an earlier conceptual
precursor~\cite{DuttRokhlin1993}; Harvey and van der Hoeven discuss the
connection in~\cite[Section~4.4.3]{HarveyHoeven2021}. The quantitative
resampling and bit-cost estimates used here are those of the latter paper.

More specifically, the proof adapts Harvey and van der Hoeven's
Gaussian-resampling identity, contraction bounds, and short-convolution
procedures~\cite[Theorem~4.2, Proposition~4.7 and Lemmas~4.8--4.12]{HarveyHoeven2021}.
Their framework already permits arbitrary dimension. Here the source and
target permutations remain in the transform identity, while the faster tape
procedures perform the required axis movements. Bluestein's conversion to
convolution~\cite{Bluestein1970} and the synthetic-root convolution algorithm
are retained; the latter follows the polynomial-transform construction
in~\cite[Section~2.4 and Section~3]{HarveyHoeven2021}. Signed coefficient
packing is a form of Kronecker substitution, with the fixed-tape version
of~\cite[Lemma~2.5]{HarveyHoeven2021} as a direct antecedent.
Calls to the established $O(n\log n)$ multiplier are ordinary subroutine
calls; they do not invoke the improved theorem recursively.

\subsection{The finite-network ingredients}

The subset-intersection construction in Section~\ref{sec:finite-networks}
has an earlier source outside multiplication. Alon's presentation of the
Frankl--Wilson construction~\cite{FranklWilson1981} gives low-degree
representations of complementary intersection graphs over different fields
\cite[Section~3, remark after the proof of Theorem~1.1]{Alon1998}.
At the prime $p=2$, its subsets have size three and the two polynomial evaluations
are the intersection count modulo two and the intersection count minus one.
These are the two pairing formulas used here. We change the ground-set
size and express the second pairing through a rational bilinear form.

The gather and scatter operations, together with the side wires, have the
form of linear index-code decoding: a low-rank linear summary gives the
wanted coordinate plus contributions from known neighboring coordinates,
and those contributions are subtracted. Bar-Yossef, Birk, Jayram and Kol
establish the fitting-matrix formulation and this decoder
\cite[Section~3, proof of Theorem~5]{BarYossefBirkJayramKol2006}.
Lubetzky and Stav treat decoding over other fields and use related
intersection representations to separate their possible ranks
\cite[Proposition~2.1 and Section~2.3]{LubetzkyStav2009}.
These comparisons explain the algebraic role of the two fields here;
they do not supply the subspace labels on the subsequent network.

Restoring scratch registers with arbitrary initial contents also has a
precise precedent. The transparent-computation construction of Buhrman,
Cleve, Kouck\'y, Loff and Speelman subtracts an initial register value,
performs a computation, adds the resulting value, and reverses the
computation~\cite[Section~3]{BuhrmanCleveKouckyLoffSpeelman2014}.
Our scratch schedule is the linear-readout version of that construction.
The rank reduction used for address shears is a singular-matrix Bruhat
factorization: after reversing the coordinate order, Grigor'ev's
upper-triangular factors become the two lower-triangular factors needed
here~\cite[Appendix~A.2, Proposition~14(c)]{Grigoriev1982}.

The finite networks and their labels are proved directly in
Section~\ref{sec:finite-networks}. The additional requirements are the
three-stage bank exchange, nested subspaces with a strict total rank or
dimension deficit, and endpoint frames acting on every role. The tape
implementations then turn that deficit into the volume-normalized
recurrences of Sections~\ref{sec:tape-interchange}
and~\ref{sec:fast-butterfly-layers}, including the costs of addressing,
padding, coefficient growth, and finite precision. The cited decoding,
restoration, and factorization results provide ingredients; these further
properties are established in the present proof.

\subsection{Lower-bound scope and matrix transposition}

Sch\"onhage and Strassen conjectured $n\log n$ as the optimal order of
multiplication cost~\cite[Section~1]{SchonhageStrassen1971}. In the present
fixed-tape model, an eventual lower bound for every multiplication machine means that, for
each fixed machine $A$, there are constants $c_A>0$ and $N_A$ such that
$T_A(n)\ge c_A n\lg n$ whenever $n\ge N_A$.
Its negation requires one machine satisfying
$\liminf_{n\to\infty}T_A(n)/(n\lg n)=0$.
The power saving proved here gives a limit of zero through all input lengths,
which is a stronger conclusion.

Harvey and van der Hoeven consider a binary transposition machine
supplied an integer $m\ge1$ and a $k\times k$ binary matrix, where
$k=\lfloor\sqrt m\rfloor$. They proved that if every fixed such machine
has worst-case cost $\Omega(m\lg m)$ through all sufficiently large
supplied values of $m$, then every fixed multiplication machine has the
corresponding eventual lower bound~\cite[Theorem~1.2]{HarveyHoeven2025}.
Combining their quantitative reduction~\cite[Corollary~6.2]{HarveyHoeven2025}
with Theorem~\ref{thm:main} gives a deterministic fixed-tape algorithm
that transposes a row-major $k\times k$ binary matrix in
$O(k^2(\lg k)^{1-\kappa})$ time, with $\kappa$ as in
Theorem~\ref{thm:main}. Section~\ref{sec:transposition} proves the more
general rectangular-matrix bound.

This comparison concerns unrestricted computation on a fixed number of
one-dimensional tapes. Our address procedure forms XOR combinations of
intermediate data, even though its final effect is a permutation.
It therefore lies outside models restricted to moving data without such
computations; see~\cite[Section~1.4]{HarveyHoeven2025} for the distinction.

\subsection{Companion application to exact Fourier transforms}

The complex phase network of Proposition~\ref{prop:complex-motif-interface}
is also used in the companion manuscript~\cite[Section~1.1, Theorem~1.1
and Corollary~1.2]{OpenAIExactDFT2026}. That manuscript gives a single
algorithm for the exact discrete Fourier transform at every positive
transform length $n$. As $n\to\infty$, its cost is
$O(n(\log n)^{1-10^{-13}})=o(n\log n)$.

The cost is measured in an exact complex-register model that assigns unit
cost to exact complex field operations. The algorithm accepts every
$x\in\mathbb C^n$, with no bound on its entries or on intermediate complex
magnitudes. Its computation on the input uses only addition, subtraction,
and multiplication by prepared input-independent scalars. The magnitudes
of those scalar coefficients are unrestricted as well.

Integer operations and random access on $O(\log(n+2))$-bit words have unit
cost, with a fixed number of words for each address or integer. The model
supplies one specified Fourier root $\zeta_{D_*}$ of explicitly computable
order $D_*<1024n^3$. Scalar preparation, schedule construction, and index
work are charged. This is a separate application of the finite network,
with its own uniform array implementation. It is neither a bit-complexity
consequence nor a stable floating-point FFT consequence of the fixed-tape
multiplication theorem, Theorem~\ref{thm:main}.
